% GENERATED FILE. Edit canonical lessons, then run npm run generate:books.
% canonical-source-hash: fnv1a64:d8cd71f313d1d237
% canonical-lessons: HI-C35-lena, HI-C35-puchna, HI-C35-madad, HI-C35-pasand

\chapter{Taking, Asking, Helping — and Liking}
\label{ch:hi-pasand}
\hlchaptermodality{35}

\begin{chapteropening}
\textbf{By the end of this chapter:} \emph{I can say I take, ask and help in Hindi, build a conjunct verb out of any noun plus \hi{करना}, and say what I like the way Hindi says it — with the thing liked as the subject and me in the dative.}

Say \hi{लेना}, \hi{पूछना} and \hi{मदद} \hi{करना}, then turn the sentence inside out: \hi{मुझे} \hi{रोटी} \hi{पसंद} \hi{है} is bread being pleasing to me, so \hi{है} agrees with the bread and I move into \hi{मुझे} — the same inversion Spanish and Italian reached on their own.
\end{chapteropening}

\section[lenā]{\hi{लेना} (lenā) — \textquotedblleft{}to take,\textquotedblright{} a verb that lost a consonant}

\label{lesson:HI-C35-lena}



\begin{quote}
Four verbs of the mind are behind you. This chapter puts your hands back in the world — starting with the shortest verb in it.
\end{quote}

\begin{cousinweb}
\textbf{\hi{लेना}} (\emph{lenā}) = \textquotedblleft{}\textbf{to take}.\textquotedblright{} Two syllables, and one of the most-used verbs in the language. Stem \textbf{\hi{ले}-}: \textbf{\hi{मैं} \hi{लेता} \hi{हूँ}} (m.) / \textbf{\hi{मैं} \hi{लेती} \hi{हूँ}} (f.); \textbf{\hi{आप} \hi{लेते} \hi{हैं}}.

It reaches further than English \textquotedblleft{}take\textquotedblright{}. Offered a drink, a Hindi speaker answers \textbf{\hi{मैं} \hi{पानी} \hi{लेता} \hi{हूँ}} — \textquotedblleft{}I'll have water,\textquotedblright{} the same \emph{lenā}. It pairs with its opposite in the fixed phrase \textbf{\hi{लेन}-\hi{देन}} (\emph{len-den}), \textquotedblleft{}\textbf{taking and giving}\textquotedblright{} — Hindi's word for dealings, transactions, business between people.
\end{cousinweb}

\subsection*{You'll want to know — the consonant that fell out}
\textbf{\hi{लेना}} looks too short to have a history. It has a long one. Sanskrit's verb was \textbf{\hi{लभते}} (\emph{labhate}) — \textquotedblleft{}\textbf{takes, obtains, receives}.\textquotedblright{} Through Prakrit \emph{lahaï} the \textbf{-bh-} thinned to a breath and then vanished, and the two vowels left behind fused into the single \textbf{\hi{ले}-} you now say. Almost the whole word eroded.

The proof is sitting in Hindi's own vocabulary. \textbf{\hi{लाभ}} (\emph{lābh}) — \textquotedblleft{}\textbf{profit, gain, benefit}\textquotedblright{} — was borrowed back from Sanskrit later, straight out of the books, so it never went through that erosion. Verb and noun are the same word at two different ages: one worn smooth by mouths, one preserved on the shelf.

Further out, the standard comparison lines \emph{labh-} up with Greek \textbf{lambánein}, \textquotedblleft{}to take\textquotedblright{} — and English quietly holds three of its descendants: a \textbf{syllable} is letters \textquotedblleft{}taken together\textquotedblright{}; \textbf{epilepsy} is \textquotedblleft{}a seizing upon\textquotedblright{}; so is \textbf{catalepsy}.

\subsection*{Your turn}
\begin{itemize}
\raggedright
  \item {[}YOU SAY (m.): \emph{maiṁ letā hūṁ}  ·  (f.): \emph{maiṁ letī hūṁ}{]}
  \item \emph{Say it:} accept a drink — \emph{maiṁ pānī letā hūṁ}; then a meal — \emph{maiṁ roṭī letā hūṁ}. Object first, verb last, as always
  \item \emph{Say it:} \emph{pānī} and what it literally means (the drinkable thing), then \emph{roṭī}
  \item \emph{Say it:} \emph{lenā} and \emph{lābh} one after the other — the same Sanskrit verb, worn down and preserved
\end{itemize}

\begin{tcolorbox}[breakable,colback=teal!4,colframe=teal!35!black,title={Before you move on}]
Which Sanskrit verb is \emph{lenā}, and which consonant did it lose? (\textbf{\hi{लभते}} \emph{labhate}; the \textbf{-bh-}.) Which Hindi noun still shows that consonant, and why? (\textbf{\hi{लाभ}} — borrowed back from Sanskrit, so it never eroded.) Say \textquotedblleft{}I'll have water.\textquotedblright{} (\emph{Maiṁ pānī letā / letī hūṁ}.)
\end{tcolorbox}
\section[\hi{पूछना}]{\hi{पूछना} (pūchnā) — \textquotedblleft{}to ask,\textquotedblright{} and the reason English says \textquotedblleft{}pray\textquotedblright{}}

\label{lesson:HI-C35-puchna}



\begin{quote}
You have been asking questions since chapter two. Here, at last, is the verb for doing it.
\end{quote}

\subsection*{The letters in this word}
\textbf{\hi{प}} (\emph{pa}) with the long \textbf{\hi{ू}} hanging beneath it $\to$ \textbf{\hi{पू}} (\emph{pū}); then \textbf{\hi{छ}} (\emph{cha}), the breathy pair of \emph{ca}; then \textbf{\hi{ना}} (\emph{nā}) $\to$ \textbf{\hi{पू}·\hi{छ}·\hi{ना}}.

The \textbf{ū} is genuinely long, and \textbf{\hi{छ}} takes a puff of air, one degree heavier than the \textbf{\hi{च}} in \emph{sochnā}.

\begin{cousinweb}
\textbf{\hi{पूछना}} (\emph{pūchnā}) = \textquotedblleft{}\textbf{to ask},\textquotedblright{} in the sense of asking a \textbf{question}. Stem \textbf{\hi{पूछ}-}: \textbf{\hi{मैं} \hi{पूछता} \hi{हूँ}} (m.) / \textbf{\hi{मैं} \hi{पूछती} \hi{हूँ}} (f.); \textbf{\hi{आप} \hi{पूछते} \hi{हैं}}.

You already own a question — \textbf{\hi{कितने} \hi{साल} \hi{के} \hi{हो}}, \textquotedblleft{}how many years are you of.\textquotedblright{} Now you can name the act itself: \textbf{\hi{मैं} \hi{उम्र} \hi{पूछता} \hi{हूँ}}, \textquotedblleft{}I ask (someone's) age.\textquotedblright{} The thing asked about sits in front of the verb, where every object has gone.
\end{cousinweb}

\subsection*{You'll want to know — asking a person, asking a god}
Behind \textbf{\hi{पूछना}} stands Sanskrit \textbf{\hi{पृच्छति}} (\emph{pṛcchati}, \textquotedblleft{}asks\textquotedblright{}), and behind that a reconstructed root \emph{\textbf{prek-}}, \textquotedblleft{}\textbf{to ask}.\textquotedblright{} This one paid out generously across the family:

\begin{itemize}
\raggedright
  \item Latin \textbf{precārī}, \textquotedblleft{}to beg, to entreat\textquotedblright{} — English \textbf{pray}, \textbf{prayer}, \textbf{deprecate}, \textbf{imprecation}.
  \item Latin \textbf{precārius}, \textquotedblleft{}held only so long as another allows\textquotedblright{} — English \textbf{precarious}. A precarious thing is one you keep by asking.
  \item German \textbf{fragen}, Dutch \textbf{vragen} — everyday \textquotedblleft{}to ask\textquotedblright{} to this day.
\end{itemize}

English is the odd one out, instructively. Its own inherited cousin, Old English \textbf{frignan}, died out, and \textbf{ask} comes from somewhere else entirely — so the \emph{prek-} word survived in English only in its religious clothes. Saying \textbf{\hi{मैं} \hi{पूछता} \hi{हूँ}}, you use the verb English kept for \textbf{praying}, because asking a person and asking a god were once one act with one name.

\subsection*{Your turn}
\begin{itemize}
\raggedright
  \item {[}YOU SAY (m.): \emph{maiṁ pūchtā hūṁ}  ·  (f.): \emph{maiṁ pūchtī hūṁ}{]}
  \item \emph{Say it:} the chapter-19 question, then name what you just did — \emph{kitne sāl ke ho?} … \emph{maiṁ umr pūchtā hūṁ}
  \item \emph{Say it:} \emph{maiṁ pūchtā hūṁ}, then \emph{maiṁ letā hūṁ} — I ask, I take
  \item \emph{Say it:} \emph{pūchnā}, then \textbf{pray}, \textbf{precarious}, \textbf{fragen} — one root, three destinations
\end{itemize}

\begin{tcolorbox}[breakable,colback=teal!4,colframe=teal!35!black,title={Before you move on}]
Which English verb is \emph{pūchnā}'s true cousin — \emph{ask} or \emph{pray}? (\textbf{Pray} — \emph{ask} is unrelated.) What does \emph{precarious} literally describe? (Something \textbf{held only by asking}.) Say \textquotedblleft{}I ask (someone's) age.\textquotedblright{} (\emph{Maiṁ umr pūchtā / pūchtī hūṁ}.)
\end{tcolorbox}
\section[\hi{मदद} \hi{करना}]{\hi{मदद} \hi{करना} (madad karnā) — \textquotedblleft{}to help,\textquotedblright{} and the machine that builds it}

\label{lesson:HI-C35-madad}



\begin{quote}
This one is not a verb. It is a noun with a verb attached — and once you spot the join, you can build hundreds more.
\end{quote}

\begin{cousinweb}
\textbf{\hi{मदद}} (\emph{madad}) is a \textbf{noun}: \textquotedblleft{}\textbf{help, aid}.\textquotedblright{} On its own it does nothing. Bolt \textbf{\hi{करना}} (\emph{karnā}, \textquotedblleft{}to do\textquotedblright{}) on the end and you have the verb: \textbf{\hi{मदद} \hi{करना}}, literally \textquotedblleft{}\textbf{to do help}.\textquotedblright{}

Only \textbf{\hi{करना}} ever conjugates. The noun sits still:

\begin{itemize}
\raggedright
  \item \textbf{\hi{मैं} \hi{मदद} \hi{करता} \hi{हूँ}} (m.) / \textbf{\hi{मैं} \hi{मदद} \hi{करती} \hi{हूँ}} (f.) — \textquotedblleft{}I help.\textquotedblright{}
  \item \textbf{\hi{आप} \hi{मदद} \hi{करते} \hi{हैं}} — \textquotedblleft{}you (resp.) help.\textquotedblright{}
\end{itemize}

This is the \textbf{conjunct verb}, and you have already used two without being told: \textbf{\hi{काम} \hi{करना}}, \textquotedblleft{}to do work\textquotedblright{} = to work; and \textbf{\hi{माफ़} \hi{कीजिए}}, \textquotedblleft{}please do forgiveness\textquotedblright{} = sorry. One join, three verbs.
\end{cousinweb}

\subsection*{You'll want to know — a stretched-out hand, and why the pattern matters}
\textbf{\hi{मदद}} is Arabic \textbf{madad}, on the root \textbf{\ar{م}-\ar{د}-\ar{د}} (\emph{m-d-d}), \textquotedblleft{}\textbf{to stretch out, to extend}.\textquotedblright{} Help, in Arabic, is a thing \textbf{extended} toward you — the picture English reaches for when it \emph{extends} a hand. It came into Hindi through \textbf{Persian}, the same road as \textbf{\hi{उम्र}} (age) and \textbf{\hi{माफ़}} (forgiven).

State the gap plainly: English inherited \textbf{nothing} from \emph{m-d-d}, and English \textbf{help} is a Germanic word unrelated to any of this. The picture is shared; the word is not.

Now why this matters more than one verb. A borrowed noun cannot be conjugated — Hindi would have to bend a foreign word through its own endings. With \textbf{\hi{करना}} it never has to: the loan stays a frozen noun and native grammar does the work. That is how Hindi swallowed centuries of Persian and Arabic vocabulary without disturbing its verb system at all.

\subsection*{Your turn}
\begin{itemize}
\raggedright
  \item {[}YOU SAY (m.): \emph{maiṁ madad kartā hūṁ}  ·  (f.): \emph{maiṁ madad kartī hūṁ}{]}
  \item \emph{Say it:} the three joins in a row — \emph{kām karnā}, \emph{māf karnā}, \emph{madad karnā}; hear that only the second word ever moves
  \item \emph{Say it:} \emph{maiṁ pūchtā hūṁ}, then \emph{maiṁ madad kartā hūṁ} — I ask, I help
  \item \emph{Say it:} what \emph{m-d-d} means underneath \emph{madad} (to stretch out, extend)
\end{itemize}

\begin{tcolorbox}[breakable,colback=teal!4,colframe=teal!35!black,title={Before you move on}]
Which half of \emph{madad karnā} takes the endings? (\textbf{\hi{करना}} — the noun never changes.) Name two conjunct verbs you had already used. (\textbf{\hi{काम} \hi{करना}}; \textbf{\hi{माफ़} \hi{कीजिए}}.) Why does this pattern let Hindi borrow so freely? (The loan \textbf{stays a noun}; native \emph{karnā} carries the grammar.)
\end{tcolorbox}
\section[\hi{पसंद}]{\hi{पसंद} (pasand) — Hindi does not like things; things please Hindi}

\label{lesson:HI-C35-pasand}



\begin{quote}
Every verb so far let you be the one doing something. This last one takes that away.
\end{quote}

\subsection*{You'll want to know — \hi{मुझे}, the shape of \textquotedblleft{}to me\textquotedblright{}}
\textbf{\hi{मुझे}} (\emph{mujhe}) = \textquotedblleft{}\textbf{to me}.\textquotedblright{} It is what \textbf{\hi{मैं}} (\textquotedblleft{}I\textquotedblright{}) becomes when the sentence will not let you be its subject. \emph{Mujhe} is not \textquotedblleft{}I.\textquotedblright{}

\begin{grammarlens}[title={Grammar Lens: the thing you like is the subject}]
\textbf{\hi{मुझे} \hi{रोटी} \hi{पसंद} \hi{है}} = \textquotedblleft{}\textbf{I like bread}\textquotedblright{} — or, taken apart: \textbf{\hi{मुझे}} (\textquotedblleft{}to me\textquotedblright{}) + \textbf{\hi{रोटी}} (\textquotedblleft{}bread\textquotedblright{}) + \textbf{\hi{पसंद}} (\textquotedblleft{}pleasing\textquotedblright{}) + \textbf{\hi{है}} (\textquotedblleft{}is\textquotedblright{}). \textquotedblleft{}\textbf{To me, bread is pleasing.}\textquotedblright{}

\textbf{\hi{रोटी}} is the \textbf{subject}, and that is what \textbf{\hi{है}} agrees with. You are only the person it happens to. Hindi has no \textquotedblleft{}I like bread\textquotedblright{} with \emph{I} in front.

Two languages walked here alone: Spanish \textbf{me gusta el café} (\textquotedblleft{}to me, the coffee pleases\textquotedblright{}) and Italian \textbf{mi piace}. Those two are close cousins; Hindi is not, and nobody borrowed it. Two families, one decision — liking happens \textbf{to} you.
\end{grammarlens}

\subsection*{You'll want to know — a Persian word doing a verb's job}
Because \textbf{\hi{पसंद} is not a verb}. It is \textbf{Persian} — from \textbf{\ar{پسندیدن}} (\emph{pasandīdan}, \textquotedblleft{}to approve, to find pleasing\textquotedblright{}) — and it entered Hindi as a noun and adjective, \textquotedblleft{}\textbf{liked, one's choice}.\textquotedblright{} A non-verb cannot carry the person doing the liking, so the grammar put them elsewhere: \textbf{\hi{मुझे}}.

It came by the Persian road that also gave you \textbf{\hi{शाम}} (evening) — and see what these two chapters have shown. Hindi's verbs of the mind (\emph{sochnā, samajhnā, paṛhnā, likhnā}) are inherited Indo-Aryan to a word, while \emph{madad} and \emph{pasand} — helping and liking — are borrowed. Thinking stayed home; the social vocabulary travelled.

For an active \textquotedblleft{}like\textquotedblright{}, Hindi uses last lesson's join: \textbf{\hi{पसंद} \hi{करना}}, \textquotedblleft{}to do liking,\textquotedblright{} nearer to \textquotedblleft{}I choose, I prefer.\textquotedblright{} For everyday liking, \textbf{\hi{मुझे} … \hi{पसंद} \hi{है}}.

\subsection*{Your turn}
Say these aloud:

\begin{itemize}
\raggedright
  \item \emph{mujhe roṭī pasand hai}, then the Hindi meaning — \textquotedblleft{}to me, bread is pleasing\textquotedblright{}
  \item swap the thing, leaving \emph{mujhe … pasand hai} untouched — \emph{mujhe pānī pasand hai} · \emph{mujhe kuttā pasand hai} · \emph{mujhe shām pasand hai}
  \item the chapter end to end — \emph{maiṁ letā hūṁ}, \emph{maiṁ pūchtā hūṁ}, \emph{maiṁ madad kartā hūṁ}, \emph{mujhe roṭī pasand hai}
  \item \emph{pasand karnā} beside \emph{madad karnā} — one join, twice
\end{itemize}

\begin{tcolorbox}[breakable,colback=teal!4,colframe=teal!35!black,title={Before you move on}]
In \emph{mujhe roṭī pasand hai}, what is the subject? (\textbf{\hi{रोटी}} — which is why it is \emph{hai}.) What part of speech is \emph{pasand}, and from where? (A \textbf{noun/adjective}, from \textbf{Persian} — not a verb.) Say \textquotedblleft{}I like the dog,\textquotedblright{} then what it literally means. (\emph{Mujhe kuttā pasand hai} — \textquotedblleft{}to me, dog is pleasing.\textquotedblright{})
\end{tcolorbox}
